\begin{figure*}[t]
  \centering
  \includegraphics[width=\textwidth]{figures/overview-standalone.pdf}
  \caption{\textbf{Overview of \sys.}
  (1) \emph{Memory construction}: dialogue histories are turned into
  dated propositions, stored both as statements and as graph triples
  linked to their source utterances. (2) \emph{Memory retrieval}: a
  planner grounded in retrieved evidence specifies the relational
  requirements of the question, its answer type, and a temporal reference.
  An executor searches the graph for witnesses under consistent bindings;
  a verifier checks candidates against their sources, and gaps or
  rejections are fed back to the planner. Accepted witnesses guide the
  selection of propositions and passage summaries passed to the reader.
  Blue outlines denote LLM operations, green outlines denote retrieval or
  deterministic processing, dotted links denote provenance, and dashed
  arrows denote feedback.}
  \Description{The architecture has two stages: memory construction and memory retrieval. Memory construction
  extracts propositions and indexes their time and importance metadata.
  Source passages and the dated proposition graph are connected by
  provenance. A question triggers evidence retrieval and an LLM planner.
  The planner specifies conjunctive requirements, answer type, temporal
  reference and weights. A graph executor finds consistent bindings without
  LLM calls. A conditional verifier checks source utterances, returning
  gaps or rejections to the planner. Accepted support and relevance-ranked
  candidates guide context selection. The reader receives dated statements
  and passage summaries and generates an answer.}
  \label{fig:overview}
\end{figure*}
